\chapter{Theoretical Background}
\label{cap:fundamentacao}

This chapter introduces the concepts on which the remainder of this work is
built. Section~\ref{sec:fund-regressao} formalizes the regression problem and
the metrics used to evaluate regression models. Section~\ref{sec:fund-algoritmos}
describes the five learning algorithms adopted in the experiments.
Section~\ref{sec:fund-pipeline} presents the pre-processing and model selection
steps that compose a typical ML pipeline, together with the validation
procedures used to estimate generalization performance.
Section~\ref{sec:fund-vazamento} discusses how data leakage arises inside these
pipelines. Finally, Section~\ref{sec:fund-estatistica} summarizes the statistical
tools used to analyze paired performance estimates.

\section{Regression and Its Evaluation}
\label{sec:fund-regressao}

In supervised learning, a model is induced from a set of $n$ labeled examples
$\mathcal{D} = \{(\mathbf{x}_i, y_i)\}_{i=1}^{n}$, where each
$\mathbf{x}_i \in \mathbb{R}^{d}$ is a vector of $d$ input features and $y_i$ is
the value of the target variable. When the target is categorical the task is
called classification; when it is continuous, $y_i \in \mathbb{R}$, the task is
called regression \cite{hastie2009}. The goal of a learning algorithm is to
produce a function $\hat{f}: \mathbb{R}^d \to \mathbb{R}$ whose predictions
$\hat{y} = \hat{f}(\mathbf{x})$ are close to the true values for \emph{new}
examples drawn from the same distribution as $\mathcal{D}$. The quantity of
interest is therefore the \emph{generalization error} of $\hat{f}$, that is, its
expected loss on unseen data, and not the error measured on the examples used to
fit it.

Since the target is continuous, the quality of a regression model is measured by
how far its predictions deviate from the observed values. Given a set of $m$
evaluation examples with observed targets $y_i$ and predictions $\hat{y}_i$, the
four metrics adopted in this work are defined as follows:

\begin{align}
  \text{MAE}  &= \frac{1}{m}\sum_{i=1}^{m} \lvert y_i - \hat{y}_i \rvert, \label{eq:mae}\\
  \text{MSE}  &= \frac{1}{m}\sum_{i=1}^{m} (y_i - \hat{y}_i)^2, \label{eq:mse}\\
  \text{RMSE} &= \sqrt{\text{MSE}}, \label{eq:rmse}\\
  R^2         &= 1 - \frac{\sum_{i=1}^{m} (y_i - \hat{y}_i)^2}{\sum_{i=1}^{m} (y_i - \bar{y})^2}, \label{eq:r2}
\end{align}

\noindent where $\bar{y}$ is the mean of the observed targets. The Mean Absolute
Error (MAE) weights all residuals linearly and is therefore less sensitive to
large errors than the Mean Squared Error (MSE), which penalizes them
quadratically. The Root Mean Squared Error (RMSE) brings the MSE back to the unit
of the target variable, which makes it easier to interpret. These three metrics
are unbounded and scale-dependent: an RMSE of 10 may be excellent for a target
measured in thousands and poor for a target measured in units. As a consequence,
their absolute values cannot be compared across data sets.

The coefficient of determination ($R^2$) does not suffer from this limitation.
It expresses the fraction of the variance of the target that is explained by the
model, relative to a baseline that always predicts the mean. A perfect model has
$R^2 = 1$, a model that is as good as the mean predictor has $R^2 = 0$, and a
model that is worse than the mean predictor has $R^2 < 0$, which may happen on
held-out data. \citeonline{chicco2021} argue that, because it is normalized and
easy to interpret, $R^2$ is more informative than MAE, MSE, and RMSE when
regression models are compared across problems. For this reason, $R^2$ is the
primary metric of this study, while RMSE is reported as a complementary
measure expressed in relative terms (Chapter~\ref{cap:metodologia}).

\section{Regression Algorithms}
\label{sec:fund-algoritmos}

The algorithms considered in this work were chosen to cover distinct learning
paradigms: a parametric linear model, an instance-based method, a single tree, an
ensemble of trees, and a kernel-based method. Their different inductive biases
make them react differently to changes in feature scale, in the set of available
features, and in the choice of hyperparameters, which is precisely what the
leakage scenarios manipulate.

\subsection{Linear Regression}

Linear Regression (LR) assumes that the target is a linear function of the
features, $\hat{y} = \beta_0 + \sum_{j=1}^{d} \beta_j x_j$, and estimates the
coefficients by Ordinary Least Squares, i.e., by minimizing the sum of squared
residuals over the training set \cite{hastie2009}. It has no hyperparameters in
its basic form and its solution has a closed form. An important property for
this work is that the predictions of an OLS model with intercept are invariant to
affine transformations of the features: rescaling or shifting a feature changes
the corresponding coefficient, but not the fitted values. Optional constraints,
such as forcing all coefficients to be non-negative, turn the problem into a
constrained least-squares problem and can be treated as hyperparameters.

\subsection{K-Nearest Neighbors}

The K-Nearest Neighbors (KNN) regressor is a non-parametric, instance-based
method: it stores the training set and predicts the target of a new example as
the average target of the $k$ training examples closest to it
\cite{cover1967,hastie2009}. Its behavior is determined by the number of
neighbors $k$, which controls the bias-variance trade-off, and by the distance
metric used to define proximity (typically Euclidean or Manhattan). Because
distances are computed directly on the feature values, features with larger
ranges dominate the neighborhood structure, which makes KNN highly sensitive to
feature scaling.

\subsection{Decision Trees}

A regression Decision Tree (DT) recursively partitions the feature space with
axis-aligned splits of the form $x_j \leq t$, choosing at each node the feature
and threshold that most reduce the squared error of the resulting child nodes,
and predicts the mean target of the training examples that reach each leaf
\cite{breiman1984}. Trees capture non-linear relations and interactions without
requiring any transformation of the inputs, and the partitions they induce are
invariant to monotonic transformations of each feature. Left unconstrained, a
tree grows until its leaves are pure and overfits the training data; its
complexity is therefore controlled by hyperparameters such as the maximum depth.

\subsection{Random Forests}

Random Forests (RF) combine many decision trees, each trained on a bootstrap
sample of the training set and restricted to a random subset of candidate
features at each split, and average their predictions \cite{breiman2001}. The
randomization decorrelates the trees, so that averaging reduces the variance of
the ensemble with respect to that of a single tree. RF inherits the insensitivity
of trees to feature scaling, and its main hyperparameters are the number of trees
and the maximum depth of each tree.

\subsection{Support Vector Regression}

Support Vector Regression (SVR) adapts the maximum-margin principle of Support
Vector Machines to continuous targets \cite{drucker1997,smola2004}. It searches
for a function that deviates from the observed targets by at most $\varepsilon$
for as many training examples as possible while remaining as flat as possible,
penalizing deviations larger than $\varepsilon$ linearly. The trade-off between
flatness and tolerance to violations is controlled by the regularization
parameter $C$. Through the use of kernel functions, such as the linear and the
Radial Basis Function (RBF) kernels, SVR can model non-linear relations. Since
both the kernel values and the $\varepsilon$-insensitive loss depend on the
scale of the data, SVR is known to require standardized inputs, and often a
standardized target as well, to work properly.

\section{Building and Evaluating ML Pipelines}
\label{sec:fund-pipeline}

In practice, a learning algorithm is rarely applied to raw data. It is preceded
by a sequence of data transformations and followed by procedures that choose its
configuration and estimate its performance. This whole sequence is referred to
as an ML \emph{pipeline}. The four steps investigated in this work are described
below, followed by the validation procedures that turn a pipeline into a
performance estimate.

\subsection{Feature Scaling}

Features measured in different units may have ranges that differ by several
orders of magnitude. Scaling methods map them to comparable ranges. The most
common one, standardization (or $z$-score normalization), transforms each
feature $x_j$ into
\begin{equation}
  x'_j = \frac{x_j - \mu_j}{\sigma_j},
  \label{eq:zscore}
\end{equation}
where $\mu_j$ and $\sigma_j$ are the mean and the standard deviation of the
feature. Equation~\ref{eq:zscore} hides an important detail: $\mu_j$ and
$\sigma_j$ are \emph{estimated from data}. The scaler is therefore a model in
itself, with parameters that are learned during a fitting phase and then applied
to new data in a transformation phase. Which examples are used in the fitting
phase is exactly what distinguishes a correct from a leaky pipeline.

\subsection{Missing Value Imputation}

Real data sets frequently contain missing entries, and most learning algorithms
cannot handle them directly. The statistical literature classifies missingness
into three mechanisms \cite{rubin1976,little2019}: values are \emph{missing
completely at random} (MCAR) when the probability of a value being missing does
not depend on any data; \emph{missing at random} (MAR) when it depends only on
observed values; and \emph{missing not at random} (MNAR) when it depends on the
missing value itself. Imputation methods replace the missing entries with
plausible values. The simplest ones replace them with the mean or the median of
the observed values of the same feature. The KNN imputer replaces each missing
entry with the average of that feature over the $k$ most similar examples,
measuring similarity on the features that are observed \cite{troyanskaya2001}.
As with scaling, all imputers learn statistics from data (means, medians, or a
reference set of neighbors) before they can transform it.

\subsection{Feature Selection}

Feature selection reduces the input space to a subset of relevant features, which
may improve generalization, reduce computational cost, and make models easier to
interpret \cite{guyon2003}. Selection methods are commonly divided into filter,
wrapper, and embedded methods. Filter methods, used in this work, rank the
features according to a statistic computed independently of the learning
algorithm and keep the best-ranked ones. For regression, a standard filter is the
univariate $F$-test, which measures the strength of the linear association between
each feature and the target; one may then keep a fixed number or a fixed
percentile of the highest-scoring features. Because the ranking uses the target
values, the selection step directly consumes information about the quantity that
the model will later be asked to predict.

\subsection{Hyperparameter Tuning}

Hyperparameters are configuration values that are not learned by the training
algorithm itself, such as $k$ in KNN, the maximum depth of a tree, or $C$ and the
kernel in SVR. Tuning consists of searching a space of candidate configurations
for the one that yields the best estimated performance. Grid search, which
evaluates every combination in a predefined grid, is the most common strategy;
random search is a frequently used alternative for large spaces
\cite{bergstra2012}. Whatever the strategy, every candidate must be evaluated on
data not used to train it, so that tuning itself relies on a validation
procedure.

\subsection{Validation Procedures}

The simplest way to estimate generalization performance is the hold-out method:
the data set is split once into a training set, used to fit the pipeline, and a
test set, used only to evaluate it. Hold-out estimates have high variance on
small data sets, since they depend on a single random split. In $k$-fold
cross-validation (CV), the data are partitioned into $k$ folds of similar size;
the pipeline is trained $k$ times, each time on $k-1$ folds, and evaluated on the
remaining fold, and the $k$ scores are averaged \cite{kohavi1995,hastie2009}.
Repeating the procedure with different random partitions (repeated CV) further
reduces the variance of the estimate.

When CV is used both to choose hyperparameters and to report performance, the
reported score is the maximum over many noisy estimates and is, consequently,
optimistically biased \cite{varma2006,cawley2010}. Nested cross-validation
removes this bias by using two loops: an inner CV, run only on the training part
of each outer fold, selects the hyperparameters, and the outer CV evaluates the
selected configuration on data that took no part in the selection. Nested CV is
computationally more expensive, but it yields a nearly unbiased estimate of the
performance of the \emph{whole} tuning procedure \cite{varma2006}.

\section{Data Leakage in ML Pipelines}
\label{sec:fund-vazamento}

The validation procedures described above rest on one assumption: that the data
used for evaluation play no role whatsoever in building the model being
evaluated. Data leakage is any violation of this assumption, i.e., the use,
during model development, of information that would not be available when the
model is applied to genuinely new data \cite{kaufman2012}. Leakage may come from
the data themselves, for example when a feature encodes the target or when
duplicated records end up in both training and test sets, or from the way the
pipeline is built. \citeonline{kapoor2023} include in their taxonomy a class of
leakage caused by the lack of a clean separation between training and test data,
which covers pre-processing applied to the full data set before splitting.

This last form is the focus of this work. As discussed in
Section~\ref{sec:fund-pipeline}, scalers, imputers, feature selectors, and tuning
procedures all \emph{learn} something from the data. In a correct pipeline,
each of these components is fitted only on the training portion of each split
and then applied, frozen, to the evaluation portion. A leakage-prone pipeline
performs the same operations once, on the whole data set, and only then splits
the data. The test examples thus contribute to the statistics or decisions that
shape the model, and the evaluation is no longer performed on data that are
truly unseen. Table~\ref{tab:fontes-vazamento} summarizes, for each pipeline step
studied here, what is learned from the data and what reaches the model when the
step is fitted before splitting.

\begin{table}[htb]
\centering
\caption{What each pipeline step learns from data and what leaks when it is fitted on the full data set.}
\label{tab:fontes-vazamento}
\small
\begin{tabular}{p{2.6cm}p{5.0cm}p{6.2cm}}
\toprule
\textbf{Step} & \textbf{Learned from data} & \textbf{Information leaked from the test set} \\
\midrule
Scaling & Mean and standard deviation of each feature & Location and spread of the test feature values \\
Imputation & Mean, median, or reference neighbors of each feature & Test feature values used to fill gaps in the training set (and vice versa) \\
Feature selection & Ranking of features by their association with the target & Test \emph{target} values, through the choice of which features are kept \\
Hyperparameter tuning & Best configuration in the search space & Performance on the evaluation folds, through the choice of the configuration that is reported \\
\bottomrule
\end{tabular}
\\[2pt]
\footnotesize Source: The Author
\end{table}

The table suggests that the four mechanisms are not equally harmful. Scaling and
imputation leak information only about the distribution of the features, which,
under the usual assumption that training and test examples are drawn from the
same distribution, is already well estimated by the training data alone.
Feature selection and hyperparameter tuning, in contrast, leak information
related to the \emph{target}, either directly or through performance scores, and
can therefore adapt the model to the particular test examples used for
evaluation. These expectations are tested empirically in
Chapter~\ref{cap:resultados}.

\section{Statistical Tools for Paired Comparisons}
\label{sec:fund-estatistica}

Comparing a leakage-prone and a leakage-free pipeline on the same data splits
produces paired observations. The Wilcoxon signed-rank test \cite{wilcoxon1945}
is a non-parametric test for the null hypothesis that the differences of such
pairs are distributed symmetrically around zero. It ranks the absolute
differences and compares the sums of the ranks of positive and negative
differences, which makes it robust to outliers and does not require the
differences to be normally distributed. \citeonline{demsar2006} recommends it
for comparing two learning methods over multiple data sets, using one summary
value per data set as the unit of analysis. Zero differences are discarded by the
test, so a high proportion of exact ties reduces the effective sample size.

To compare more than two related samples, such as the effect of leakage on five
algorithms measured over the same data sets, the Friedman test \cite{friedman1937}
ranks the algorithms within each data set and tests whether their average ranks
differ more than would be expected by chance \cite{demsar2006}.

Whenever many hypotheses are tested, the probability that at least one of them is
rejected by chance grows with the number of tests. The Holm step-down procedure
\cite{holm1979} controls this family-wise error rate by sorting the $m$ $p$-values
in increasing order and comparing the $i$-th smallest with $\alpha/(m-i+1)$; it
is uniformly more powerful than the classical Bonferroni correction. Finally, the
Spearman rank correlation coefficient $\rho$ \cite{spearman1904}, a
non-parametric measure of monotonic association, is used in this work to relate
the magnitude of the leakage effect to characteristics of the data sets.
